\section{Week 10: Governance, Assurance, and the Future of AI Systems}

\subsection{Learning Objectives}

By the end of this week, you should be able to:
\begin{itemize}
  \item Reason about LLMs as socio-technical systems
  \item Identify governance and risk management needs for AI systems
  \item Construct a basic AI assurance case
  \item Distinguish technical risk from organizational and human risk
  \item Articulate limits of current LLM paradigms
\end{itemize}

\subsection{LLM Components Covered}

\textbf{System context}
\begin{itemize}
  \item Organizational embedding
  \item Human–AI interaction
  \item Lifecycle management
\end{itemize}

\textbf{Governance}
\begin{itemize}
  \item Assurance cases
  \item Risk controls
  \item Accountability mechanisms
\end{itemize}

\subsection{Conceptual Core: LLMs as Socio-Technical Systems}

LLMs do not operate in isolation.
They are embedded in:
\begin{itemize}
  \item organizations,
  \item workflows,
  \item regulatory environments,
  \item human decision-making loops.
\end{itemize}

Most real-world failures arise from \emph{system-level interactions}, not model internals.

\subsection{From Model Evaluation to System Assurance}

Benchmark scores do not constitute assurance.

Assurance requires structured arguments:
\begin{itemize}
  \item what claims are made about the system,
  \item under what assumptions,
  \item supported by what evidence.
\end{itemize}

This mirrors safety-critical engineering practices.

\subsection{Hazards and Failure Modes}

Common hazard categories:
\begin{itemize}
  \item technical (hallucination, brittleness),
  \item human (overtrust, deskilling),
  \item organizational (automation bias, diffusion of responsibility),
  \item temporal (model drift, data staleness).
\end{itemize}

Not all hazards are eliminable.

\subsection{Controls and Mitigations}

Controls operate at multiple levels:
\begin{itemize}
  \item \textbf{Preventive}: prompt constraints, retrieval grounding
  \item \textbf{Detective}: monitoring, logging, audits
  \item \textbf{Corrective}: rollback, human override, retraining
\end{itemize}

Effective systems combine technical and procedural controls.

\subsection{Assurance Cases for AI}

An assurance case typically includes:
\begin{itemize}
  \item \textbf{Claims}: what the system is safe/fit to do
  \item \textbf{Evidence}: tests, audits, evaluations
  \item \textbf{Assumptions}: operating conditions, user behavior
\end{itemize}

Assurance is never absolute — it is contextual and revisable.

\subsection{Human Factors}

Human operators:
\begin{itemize}
  \item interpret outputs,
  \item decide when to trust or override,
  \item adapt behavior over time.
\end{itemize}

Ignoring human factors is a primary cause of system failure.

\subsection{Limits of the Current Paradigm}

Current LLMs:
\begin{itemize}
  \item lack grounded understanding,
  \item have no intrinsic goals,
  \item optimize proxy objectives.
\end{itemize}

Future progress may require:
\begin{itemize}
  \item new training paradigms,
  \item tighter integration with symbolic systems,
  \item improved governance frameworks.
\end{itemize}

\subsection{Key References}

\begin{itemize}
  \item \citet{amodei2016concrete} — AI safety problems
  \item \citet{leslie2021understanding} — AI assurance
\end{itemize}

\subsection{Recommended University Lectures}

\begin{itemize}
  \item Stanford CS324 — governance and foundation models  
  \url{https://stanford-cs324.github.io/winter2022/}

  \item MIT AI Policy lectures  
  \url{https://ai.mit.edu}
\end{itemize}

\subsection{Practitioner Lab}

See \texttt{labs/week10/week10\_lab\_governance\_assurance.ipynb}.

\subsection{Week 10 Takeaway}

\begin{quote}
AI systems cannot be made safe by models alone; safety emerges from design, governance, and human judgment.
\end{quote}

